\textbf{Supplementary Materials}

\emph{Empirical constraints on the fraction of surface latent heat flux
reaching the top of atmosphere as net radiative cooling}

Ali B. Shahid

Version 5, June 2026

\textbf{S1. Forward Lagrangian trajectory analysis: methods}

Section 5.5 of the main text proposes the convective relay as the
mechanism that reconciles the low local transfer fraction (14.7 percent)
with the thermodynamic expectation that condensation heat must
eventually exit the climate system. This supplementary material reports
the direct empirical test of that mechanism using forward Lagrangian
trajectories.

\textbf{S1.1 Input data}

Trajectories were computed from the ERA5 reanalysis (Hersbach et al.,
2020) at the 200 hPa pressure level (approximately 12 km altitude,
within the upper troposphere and tropical tropopause layer). Three wind
components (zonal u, meridional v, vertical omega) were retrieved at
6-hourly resolution for a domain extending from 30 degrees S to 60
degrees N and from 90 degrees W to 30 degrees E. Seven years were
processed: 2005, 2008, 2010, 2015, 2016, 2019, and 2020, spanning
ENSO-neutral (2010, 2019), La Nina (2008, 2020), and El Nino (2015,
2016) states.

\textbf{S1.2 Release protocol}

Air parcels were released from a 1-degree grid over the Amazon basin (5
degrees S to 5 degrees N, 75 degrees W to 50 degrees W), yielding 286
parcels per release (11 latitudes x 26 longitudes). Releases were
performed on the 15th of January, April, July, and October each year at
12:00 UTC. Parcels were advected forward for 20 days using a simple
Euler integration scheme with a 6-hour timestep (80 steps total). Wind
fields at each step were interpolated to the parcel position using
bilinear interpolation. Parcels that exited the domain were retained at
their last in-domain position.

\textbf{S1.3 Control experiment}

To distinguish forest-specific transport from the general Hadley
circulation, a symmetric control experiment was conducted matching the
full main-ensemble release matrix: all 7 years (2005, 2008, 2010, 2015,
2016, 2019, 2020) and all 4 seasons (January, April, July, October
releases), for a total of 28 control runs. The control used the
identical advection protocol but released parcels from the tropical
Atlantic (5 degrees S to 5 degrees N, 30 degrees W to 5 degrees W), an
open-ocean area with no forest and no continental convective source.
The 7 years span neutral / weak La Nina (2005), La Nina (2008, 2020),
El Nino to La Nina transition (2010), strong El Nino (2015), El Nino
tail (2016), and weak El Nino (2019) phases, providing ENSO-matched
coverage of the main Amazon ensemble. If the convective relay pathway
is forest-specific, Amazon-released parcels should reach systematically
different destinations or higher-OLR zones than Atlantic-released
parcels at identical times. If the pathway reflects general Hadley
infrastructure, both sets should behave similarly.

\textbf{S1.4 CERES coupling}

For each release, the CERES EBAF Ed4.2 monthly all-sky OLR (Loeb et al.,
2018) corresponding to the release month was extracted. At each
trajectory endpoint (days 0, 5, 10, 15, 20), OLR was sampled at each
parcel location using nearest-neighbor interpolation on the 1-degree
CERES grid. The enhancement metric is the difference between the mean
OLR at parcel locations and the mean OLR over the Amazon source region
(5S-5N, 75W-50W).

\textbf{S2. Results}

\textbf{S2.1 Trajectory ensemble (Amazon source)}

Across all 28 Amazon releases, parcels consistently drift poleward and
eastward. By day 5, the mean absolute latitude reaches 16.2 degrees
(standard deviation across runs = 3.5 degrees). By day 20, the mean
latitude reaches 22.9 degrees, placing most parcels in the subtropical
Hadley descending branch. Longitudinal displacement averages 60 degrees
eastward by day 20, carrying parcels across the Atlantic into West
Africa and the eastern Atlantic.

The OLR at parcel destinations rises rapidly in the first 5 days then
plateaus. Day 5 OLR averages 270.3 W m-2 across all 28 ensembles (versus
232.1 W m-2 mean Amazon source), an enhancement of +38.2 W m-2 (standard
deviation 14.8). Day 20 OLR averages 276.7 W m-2, an enhancement of
+44.6 W m-2 (standard deviation 15.3). The enhancement is positive
across all runs and all years, confirming that parcels released from the
Amazon TTL reach higher-OLR zones within 5 days and remain there through
the 20-day integration.

\textbf{S2.2 Control experiment (Atlantic source)}

When parcels are released from the tropical Atlantic rather than the
Amazon at the identical altitude and times, the endpoint OLR values show
no consistent direction of difference from the Amazon case. Across the
full symmetric matrix of 28 control runs (7 years x 4 seasons), the
mean Amazon-minus-Atlantic endpoint OLR difference is +0.44 W m-2
(standard deviation 11.14 W m-2, range -16.4 to +23.0 W m-2). A
one-sample t-test against a null difference of zero yields t = 0.21,
p = 0.83: the null hypothesis that Amazon-source and Atlantic-source
parcels reach indistinguishable subtropical endpoint OLR cannot be
rejected. The 95\% range of differences (approximately -22 to +22 W m-2
at 2 sigma) is larger than any plausible forest-specific enhancement,
indicating that natural advection variability, not statistical
underpowering, dominates the scatter.

The result is consistent across ENSO phases. Mean Amazon-minus-Atlantic
differences by phase are: neutral / weak La Nina (2005) = -4.2 W m-2
(n = 4, standard deviation 14.0); La Nina (2008, 2020) = +0.4 W m-2
(n = 8, standard deviation 9.4); El Nino to La Nina transition (2010) =
-0.4 W m-2 (n = 4, standard deviation 13.0); strong El Nino (2015) =
+7.3 W m-2 (n = 4, standard deviation 12.7); El Nino tail (2016) =
+2.2 W m-2 (n = 4, standard deviation 10.1); weak El Nino (2019) =
-2.8 W m-2 (n = 4, standard deviation 13.2). No phase exhibits a
systematic Amazon-specific signal, and within every phase the range of
differences exceeds the phase mean by a factor of three to six.

Both Amazon-released and Atlantic-released parcels reach the subtropical
descending branch with similar endpoint OLR. The +38 to +45 W m-2
enhancement observed for Amazon parcels is therefore a general property
of tropical-to-subtropical Hadley transport, not a forest-specific
signature.

\includegraphics[width=6.25in,height=2.5in,alt={Trajectory and control analysis}]{media/media/00c08d51c3af46c6a79b903121b5ec722ca79c08.png}

\textbf{Figure S1.} Forward Lagrangian trajectory analysis from the
Amazon tropical tropopause layer at 200 hPa. (a) OLR enhancement at
parcel endpoint locations relative to the Amazon source region, averaged
across 28 trajectory ensembles (7 years, 4 seasons, 286 parcels each).
Error bars show one standard deviation. Numbers below markers indicate
mean parcel latitude. (b) Control comparison: day 20 endpoint OLR for
parcels released from the Amazon (dark green) versus the tropical
Atlantic (teal), with differences shown above each month pair. No
consistent Amazon advantage is observed.

\textbf{Table S1.} Summary of 28 Amazon trajectory ensembles. Mean and
standard deviation of parcel latitude and OLR at each endpoint day.

{\def\LTcaptype{none} % do not increment counter
\begin{longtable}[]{@{}
  >{\centering\arraybackslash}p{(\linewidth - 8\tabcolsep) * \real{0.1288}}
  >{\centering\arraybackslash}p{(\linewidth - 8\tabcolsep) * \real{0.1824}}
  >{\centering\arraybackslash}p{(\linewidth - 8\tabcolsep) * \real{0.1824}}
  >{\centering\arraybackslash}p{(\linewidth - 8\tabcolsep) * \real{0.1824}}
  >{\centering\arraybackslash}p{(\linewidth - 8\tabcolsep) * \real{0.1824}}@{}}
\toprule\noalign{}
\endhead
\bottomrule\noalign{}
\endlastfoot
\textbf{Day} & \textbf{Mean \textbar latitude\textbar{}} & \textbf{Mean
OLR (W m-2)} & \textbf{Enhancement (W m-2)} & \textbf{Standard
deviation} \\
0 & 2.7 & 233.4 & +1.3 & 1.4 \\
5 & 16.2 & 270.3 & +38.2 & 14.8 \\
10 & 20.7 & 275.0 & +42.9 & 15.4 \\
15 & 22.5 & 276.3 & +44.2 & 14.7 \\
20 & 22.9 & 276.7 & +44.6 & 15.3 \\
\end{longtable}
}

\textbf{Table S2.} Amazon versus Atlantic source control comparison, day
20 endpoint OLR for 2010. Differences have no consistent direction.

{\def\LTcaptype{none} % do not increment counter
\begin{longtable}[]{@{}
  >{\centering\arraybackslash}p{(\linewidth - 8\tabcolsep) * \real{0.1288}}
  >{\centering\arraybackslash}p{(\linewidth - 8\tabcolsep) * \real{0.1824}}
  >{\centering\arraybackslash}p{(\linewidth - 8\tabcolsep) * \real{0.1824}}
  >{\centering\arraybackslash}p{(\linewidth - 8\tabcolsep) * \real{0.1824}}
  >{\centering\arraybackslash}p{(\linewidth - 8\tabcolsep) * \real{0.1824}}@{}}
\toprule\noalign{}
\endhead
\bottomrule\noalign{}
\endlastfoot
\textbf{Month} & \textbf{Amazon OLR (W m-2)} & \textbf{Atlantic OLR (W
m-2)} & \textbf{Difference (W m-2)} & \textbf{Interpretation} \\
Jan & 273.3 & 287.9 & -14.6 & Atlantic higher \\
Apr & 273.9 & 256.4 & +17.5 & Amazon higher \\
Jul & 265.2 & 269.0 & -3.8 & Similar \\
Oct & 272.1 & 267.9 & +4.2 & Similar \\
Mean & 271.1 & 270.3 & +0.8 & No consistent signal \\
\end{longtable}
}

\textbf{S3. Interpretation}

The trajectory analysis produces two findings that together support the
main paper\textquotesingle s conclusions rather than contradict them.

First, the relay pathway is confirmed: air parcels released from the
Amazon at 200 hPa do reach subtropical latitudes within 5 to 10 days,
and those destinations have substantially higher OLR than the Amazon
source region. The mechanism described in Section 5.5 of the main text
(surface LE, cloud formation, longwave trapping, deeper convection, TTL
injection, Hadley transport, subtropical clear-sky exit) is therefore
physically real. The atmosphere does move heat from the Amazon to
subtropical clear-sky zones on weekly timescales.

Second, the pathway is not forest-specific. Parcels released from the
tropical Atlantic, where there is no forest and no continental
convective source, reach the same subtropical destinations with the same
OLR characteristics as Amazon-released parcels. The subtropical high-OLR
zone is a property of the Hadley circulation, not of forest-driven deep
convection specifically.

Taken together, these findings mean that the forest\textquotesingle s
causal contribution to TOA cooling cannot be attributed to the pathway
above 200 hPa. The pathway is shared infrastructure used by any tropical
upper-tropospheric air parcel. What the forest controls is upstream: the
quantity of heat that is injected into the TTL through deep convection.
Intact forest generates higher latent heat flux, higher convective
available potential energy, taller cumulonimbus towers, and therefore
more mass reaching the TTL than deforested land of equivalent geography.
The forest-specific TOA signal is therefore correctly quantified by the
surface-to-TOA transfer fraction at the local scale (Section 3.1) and at
the basin scale (Section 3.6), not by the trajectory endpoint OLR.

This result also constrains the spatial misattribution hypothesis,
under which local CERES measurements would underestimate forest-driven
cooling because the heat exits to space at higher latitudes. If that
were the dominant effect, Amazon parcels should reach regions with OLR
enhancement that Atlantic parcels do not. The control shows they do
not. Spatial misattribution therefore cannot account for the gap
between prior top-down estimates and the 14.7 to 20.8 percent measured
here; the geometric bound on η in deep convective regimes (main text
Section 6.10) provides the governing constraint.

\textbf{S4. Data and code availability}

Raw trajectory data (NetCDF, 28 Amazon ensembles plus 28 Atlantic
control ensembles) and the analysis scripts (Python, covering ERA5
download, trajectory integration, CERES coupling, symmetric control
comparison, and figure generation; the preloaded-vectorized control
script runs the full 56-ensemble matrix in approximately 110 seconds)
are archived with the main paper supplementary code at
\url{https://github.com/R3GENESI5/shahid-2026-transfer-fraction};
a Zenodo archive of the repository at the version of record will be
deposited on acceptance.

\textbf{References}

Hersbach, H. et al. (2020). The ERA5 global reanalysis. Quarterly
Journal of the Royal Meteorological Society, 146(730), 1999 to 2049.

Loeb, N.G. et al. (2018). Clouds and the Earth\textquotesingle s Radiant
Energy System (CERES) EBAF TOA Edition-4.0. Journal of Climate, 31(2),
895 to 918.
